\documentclass[11pt]{article}

\usepackage[margin=1.25in]{geometry}
\usepackage[T1]{fontenc}
\usepackage{amsmath,amssymb,amsthm,mathtools}
\usepackage{enumitem}
\usepackage{microtype}
\usepackage{xcolor}
\usepackage[
  colorlinks=true,
  linkcolor=blue!55!black,
  citecolor=blue!55!black,
  urlcolor=blue!55!black
]{hyperref}
\usepackage[noabbrev,nameinlink]{cleveref}

\newtheorem{theorem}{Theorem}[section]
\newtheorem{lemma}[theorem]{Lemma}
\newtheorem{proposition}[theorem]{Proposition}
\newtheorem{corollary}[theorem]{Corollary}
\theoremstyle{definition}
\newtheorem{definition}[theorem]{Definition}
\newtheorem{remark}[theorem]{Remark}

\newcommand{\qcongest}{\textsc{qCongest}}
\newcommand{\congest}{\textsc{Congest}}
\newcommand{\eps}{\varepsilon}
\newcommand{\whp}{with high probability}
\newcommand{\dist}{d_w}
\newcommand{\diamw}{\Delta_w}
\newcommand{\radw}{R_w}
\newcommand{\ecc}{e_w}
\newcommand{\polylog}{\operatorname{polylog}}
\newcommand{\wtO}{\widetilde O}
\newcommand{\wtOmega}{\widetilde\Omega}

\allowdisplaybreaks
\setlist{itemsep=2pt,topsep=4pt}

\title{Faster near-exact weighted diameter and radius in quantum CONGEST}
\author{}
\date{June 30, 2026}

\begin{document}

\maketitle

\begin{abstract}
We give randomized algorithms for weighted diameter and radius in the quantum
CONGEST model.  On a connected, undirected $n$-vertex network of unweighted
diameter $D$ and positive polynomially bounded integer edge weights, every
vertex outputs a one-sided $(1+o(1))$ overestimate, with high probability, in
\[
  \wtO\!\left(\min\{n,\,n^{5/6}+\sqrt n\,D\}\right)
\]
rounds.  Wu and Yao previously obtained
$\wtO(\min\{n^{9/10}D^{3/10},n\})$ rounds for the same problem.  For every
fixed $\delta\in(0,1/2)$, the new bound is polynomially smaller throughout
$D\leq n^{1/2-\delta}$; when $D=\Theta(\sqrt n)$, both bounds are
$\wtO(n)$.  The algorithm samples
candidate vertices independently of the landmarks used for shortest-path
reconstruction.  A candidate that is not a landmark is attached as a virtual
root to a landmark shortcut graph, and nested distributed quantum optimization
selects a good candidate and a good trial.
\end{abstract}

\begin{quote}
\small\noindent\textbf{Note.} This paper was generated by an internal harness
powered by GPT-5.5, with prompting by Aidan Bai. Cleanup and final preparation
were performed by Codex, powered by GPT-5.
\end{quote}

\section{Introduction}

Let $G=(V,E)$ be the communication network, let $w:E\to\mathbb N_{>0}$,
and write $\dist(u,v)$ for weighted shortest-path distance.  The weighted
diameter and radius are
\[
  \diamw=\max_{u,v\in V}\dist(u,v),
  \qquad
  \radw=\min_{u\in V}\max_{v\in V}\dist(u,v).
\]
We study the number of synchronous rounds needed to approximate these values
when each edge carries $O(\log n)$ qubits per round.  The parameter $D$ denotes
the diameter of the unweighted communication graph.

There is a nearly linear classical exact baseline.  Bernstein and Nanongkai's
randomized algorithm computes exact weighted all-pairs shortest paths in
$\wtO(n)$ rounds with high probability; their verification and restart
procedure gives a Las Vegas implementation with $\wtO(n)$ expected
rounds~\cite{BernsteinNanongkai2019}.  Thus exact weighted diameter and radius
also have a $\wtO(n)$ baseline.  At the other end, for every fixed
$\eps\in(0,1/2)$, Abboud, Censor-Hillel, and Khoury proved an
$\Omega(n/\log^3 n)$ classical lower bound, applying also to randomized
algorithms with constant success probability, for
$(3/2-\eps)$-approximating diameter or radius, even on sparse unweighted
networks~\cite{AbboudCensorHillelKhoury2016}.  Exact weighted single-source
shortest paths can be computed, with high probability, by the randomized
Broadcast CONGEST algorithm of Chechik and Mukhtar in
$\wtO(\sqrt n\,D^{1/4}+D)$ rounds~\cite{ChechikMukhtar2020}; its
polynomially bounded nonnegative-integer weight setting includes ours.  One
arbitrary source gives a factor-$2$ baseline for diameter and radius, but it
does not give a near-exact answer.
At a coarser guarantee, Ancona, Censor-Hillel, Dalirrooyfard, Efron, and
Vassilevska Williams give $(2+\eps)$-approximations to weighted undirected
diameter, radius, and all eccentricities with high probability in
$\wtO(\sqrt n+D)$ rounds for
$\eps=1/\polylog(n)$~\cite{AnconaEtAl2021}.  This differs in approximation
strength from the task here.

The quantum CONGEST model was introduced by Elkin, Klauck, Nanongkai, and
Pandurangan~\cite{ElkinEtAl2014}.  For unweighted networks, Le Gall and
Magniez obtained, with high probability, an exact-diameter algorithm using
$\wtO(\sqrt{nD})$ rounds and developed the distributed quantum optimization
framework behind it~\cite{LeGallMagniez2018}.  The framework of van Apeldoorn
and de Vos gives a bounded-error exact-radius algorithm at the same scale and
treats distributed quantum queries
systematically~\cite{vanApeldoornDeVos2022}.  These results do
not directly extend to near-exact weighted distances: reconstructing a
weighted shortest path requires bounded-hop distance estimates and a landmark
overlay, rather than breadth-first distances alone.  Stronger output tasks do
not share the same speedup: Wang, Wu, and Yao prove
$\Omega(n/\log n)$ quantum lower bounds for exact unweighted all-pairs
shortest paths and $\Omega(n/\log^2 n)$ for all exact eccentricities, even at
constant hop diameter~\cite{WangWuYao2021}.

Wu and Yao studied precisely this weighted problem~\cite{WuYao2022}.  Combining
distributed quantum optimization with Nanongkai's weighted shortest-path
framework~\cite{Nanongkai2014}, they obtained one-sided $(1+o(1))$
approximations with high probability in
\[
  \wtO\!\left(\min\{n^{9/10}D^{3/10},n\}\right)
\]
rounds.  They also proved that, for every fixed $\eps\in(0,1/2]$, any quantum
CONGEST algorithm that succeeds with probability at least $11/12$ and returns
a $(3/2-\eps)$-approximation requires
$\Omega(n^{2/3}/\log^2 n)$ rounds, even on graphs with
$D=\Theta(\log n)$.  Thus their lower bound leaves a polynomial gap already at
$D=\Theta(\log n)$.

Our result improves the Wu--Yao upper bound.

\begin{theorem}[Main theorem]
\label{thm:main}
Let $G$ be a connected, undirected $n$-vertex network with unweighted diameter
$D$, and suppose that every edge has a positive integer weight at most
$n^{O(1)}$.  There are randomized quantum CONGEST algorithms for weighted
diameter and radius.  With high probability, the diameter algorithm makes
every vertex output $\widehat\Delta$, and the radius algorithm makes every
vertex output $\widehat R$, satisfying
\[
  \diamw\leq \widehat\Delta\leq(1+o(1))\diamw,
  \qquad
  \radw\leq \widehat R\leq(1+o(1))\radw.
\]
Each algorithm uses
\[
  \wtO\!\left(\min\{n,\,n^{5/6}+\sqrt n\,D\}\right)
\]
rounds.  Running the two algorithms consecutively gives both values within the
same asymptotic bound.
\end{theorem}

For $D=n^\delta$ and every fixed $0\leq\delta<1/2$, the exponent in
\Cref{thm:main} is strictly smaller than the exponent in the Wu--Yao bound.
More generally, its displayed algebraic bound is $o(n)$ whenever
$D=o(\sqrt n)$, whereas the previous expression reaches its $\wtO(n)$ cap at
$D=n^{1/3}$.  There is no claimed asymptotic improvement at
$D=\Theta(\sqrt n)$, and the known $\wtOmega(n^{2/3})$ lower bound still leaves
a polynomial gap from our $\wtO(n^{5/6})$ bound for small $D$.

\paragraph{Proof idea.}
Wu and Yao use a random sample both as the domain over which an eccentricity is
optimized and as the landmark set that hits long shortest paths.  We separate
these roles.  In each of $n$ trials, a candidate set $C_i$ has expected size
$c$, while an independent landmark set $L_i$ has expected size $r$.  Every
fixed vertex appears in $\Theta(c)$ candidate sets with high probability, so
an outer quantum threshold search needs only
$\wtO(\sqrt{n/c})$ trial evaluations to encounter a diameter endpoint or a
radius center.  The landmark sample can remain smaller and is used only for
distance reconstruction.

This separation means that a candidate $s\in C_i$ need not lie in $L_i$.
Bounded-hop estimates from $s$ to the landmarks define a star from a virtual
source to the existing landmark shortcut graph.  We show that Nanongkai's
scaled bounded-hop routine can simulate these star edges by local
initialization at their landmark endpoints; no candidate-specific shortcut
graph is needed.  The path-hitting property then reconstructs every
$s$-to-$v$ distance within a factor $(1+\eps)^2$.

One trial has preprocessing and evaluation cost
\[
  \wtO\!\left(
    D+\frac nr+c+rk
    +\sqrt c\left(D+\frac{rD}{k}+r\right)
  \right),
  \tag{1}\label{eq:trial-intro}
\]
where $k$ is the shortcut parameter.  For $D\leq n^{1/3}$, the choices
$r=n^{1/3}$, $c=n^{2/3}$, and $k=\Theta(n^{1/6}\sqrt D)$ give
$\wtO(n^{5/6})$ rounds after the outer search.  For
$n^{1/3}<D\leq n^{1/2}$, the choices $r=n^{1/3}$,
$c=n^{4/3}/D^2$, and $k=\Theta(r)$ give
$\wtO(\sqrt nD)$.  We also verify that the bounded-error inner optimization
can be amplified and uncomputed to form a clean oracle for the outer quantum
search.

The output is only the value of the diameter or radius, not a witnessing pair
or center.  The algorithms are randomized and succeed with high probability;
the result is restricted to connected undirected graphs with positive
polynomially bounded integer weights.  No algorithmic or novelty claim beyond
this input and output regime is intended.

\Cref{sec:preliminaries} records the model and imported primitives.
\Cref{sec:virtual-root} proves the virtual-root reduction,
\Cref{sec:algorithm} analyzes separated sampling and coherent trial
evaluation, and \Cref{sec:outer} completes the outer search and parameter
optimization.

\section{Model and source results}
\label{sec:preliminaries}

Every vertex has a distinct $O(\log n)$-bit identifier, knows $n$, and knows
the weights of its incident edges.  Local computation and memory are
unbounded.  In one quantum CONGEST round, each endpoint of an edge may send
$O(\log n)$ qubits to the other endpoint.
Classical protocols are valid quantum protocols.  All displayed outputs have
$O(\log n)$ bits under the polynomial weight bound.  An event holds \emph{with
high probability} if its failure probability is $n^{-A}$ for an arbitrarily
large fixed constant $A$, after constants in the algorithm are increased.
The notation $\wtO$ suppresses polylogarithmic factors in $n$.

For an integer $h\geq 0$, let $\dist^h(u,v)$ be the minimum weight of a
$u$--$v$ walk using at most $h$ edges, with value $+\infty$ if no such walk
exists.  Because weights are positive, a minimum walk is a simple path.  The
eccentricity of $s$ is
\[
  \ecc(s)=\max_{v\in V}\dist(s,v),
\]
so $\diamw=\max_s\ecc(s)$ and $\radw=\min_s\ecc(s)$.

Fix one canonical shortest path $P_{uv}$ for every ordered pair $(u,v)$,
breaking ties by the sequence of vertex identifiers.  A set $L\subseteq V$ is
\emph{$h$-hitting} if every block of $h$ consecutive vertices on every
canonical path contains a vertex of $L$.  Thus, on a canonical path, the first
and last landmark are fewer than $h$ edges from the respective endpoints, and
two consecutive landmarks are at most $h$ edges apart.

We use the following consequences of Nanongkai's weighted shortest-path
framework.  Wu and Yao give the quantum-compatible formulation and account
for the pipelined broadcasts~\cite[Sec.~3 and App.~A]{WuYao2022}; in the
arXiv v2 numbering, the underlying bounded-hop and shortcut statements are
Nanongkai's Theorems 3.2, 3.6, 3.10, and 4.5 and
Lemmas 3.7 and 4.6~\cite{Nanongkai2014}.

\begin{lemma}[Bounded-hop distances and shortcuts]
\label{lem:source-shortcuts}
Let $S\subseteq V$, let $h\geq1$, and let $0<\eps\leq1$.  In
\[
  \wtO(D+h/\eps+|S|)
\]
rounds, all vertices can obtain values $\widetilde d^h(s,v)$, for
$s\in S$ and $v\in V$, such that, with high probability,
\[
  \dist(s,v)\leq \widetilde d^h(s,v)
  \leq(1+\eps)\dist^h(s,v).
  \tag{2}\label{eq:hop-estimate}
\]
Let $H$ be the graph on $m=|S|$ vertices whose finite edge weights are these
bounded-hop estimates between sources.  For every integer $1\leq k<m$, a
$k$-shortcut graph $Q$ can be constructed and embedded in
$\wtO(D+mk)$ rounds.  It satisfies
\[
  d_Q(x,y)=d_H(x,y)\quad\text{for all }x,y\in S,
  \qquad
  \operatorname{SPD}(Q)<\frac{4m}{k},
  \tag{3}\label{eq:shortcut-properties}
\]
where $\operatorname{SPD}$ is the largest number of edges on a
minimum-hop shortest path.  Every shortcut edge represents an $H$-path and
hence has weight at least the true distance between its endpoints in $G$.
\end{lemma}

We will apply a threshold form of distributed quantum optimization.  It is
important that the threshold need not be known to the algorithm.

\begin{lemma}[Distributed quantum optimization]
\label{lem:quantum-optimization}
Suppose a leader can prepare a search state on a finite domain $X$, and a
distributed unitary evaluation oracle and its inverse use $T$ rounds.  If an
unknown threshold $M$ is met by total search probability at least $\rho$, a
marked $x$ can be found, with high probability, using
$\wtO(T/\sqrt\rho)$ rounds, in addition to one-time initialization.  The same
statement holds for a lower threshold.  Exact maximum or minimum over a domain
of size $q$ uses $\wtO(\sqrt q\,T)$ rounds.
\end{lemma}

\begin{proof}
The threshold statement is Wu and Yao's Lemma 3.1, which is the unknown-threshold
form of the distributed optimization framework of Le Gall and
Magniez~\cite{WuYao2022,LeGallMagniez2018}.  Replacing values by their
order-reversing encodings gives the lower-threshold and minimum forms.  Error
reduction contributes only logarithmic factors.
\end{proof}

Finally, exact weighted all-pairs shortest paths are available in
$\wtO(n)$ rounds~\cite{BernsteinNanongkai2019}.  A BFS-tree convergecast of
the local eccentricities and a broadcast then give exact diameter or radius.
We use this as a fallback.  We also elect a leader and compute a BFS tree in
$O(D)$ rounds.  If $D_0$ is the leader's unweighted eccentricity, then
$D_0\leq D\leq2D_0$; hence the algorithms need no exact value of $D$ as input.
Using $2D_0$ in the parameter choices changes only constants, and below we
write $D$ for this estimate.

\section{A virtual root for the landmark overlay}
\label{sec:virtual-root}

Let $L$ be a landmark set of size $m$, and suppose that
$\widetilde d^h(x,v)$ is available for every $x\in L\cup\{s\}$ and
$v\in V$.  Form the landmark graph $H_L$ with edge weights
$\widetilde d^h(x,y)$, and let $Q_L$ be a $k$-shortcut graph for $H_L$.
For each landmark $x$, define
\[
  \alpha_s(x)=\widetilde d^h(s,x).
\]
An infinite value is simply omitted.  The root-initialized overlay distance is
\[
  \delta_{L,s}(u)
  =\min_{x\in L}\bigl(\alpha_s(x)+d_{Q_L}(x,u)\bigr).
  \tag{4}\label{eq:root-overlay-distance}
\]

\begin{lemma}[Virtual-root overlay]
\label{lem:virtual-root}
Let $1\leq k<m$.  From the endpoint labels $\alpha_s(x)$ and the embedded
graph $Q_L$, the network can compute and disseminate values
$\widetilde D_{L,s}(u)$ satisfying
\[
  \delta_{L,s}(u)
  \leq \widetilde D_{L,s}(u)
  \leq(1+\eps)\delta_{L,s}(u)
  \qquad(u\in L)
  \tag{5}\label{eq:virtual-root-guarantee}
\]
in
\[
  \wtO\!\left(D+\frac{mD}{\eps k}+m\right)
  \tag{6}\label{eq:virtual-root-cost}
\]
rounds, with high probability.  The construction uses no communication edge
incident to $s$ other than the original network edges and does not rebuild the
shortcut graph.
\end{lemma}

\begin{proof}
Introduce a formal vertex $\sigma$ and, for every finite $\alpha_s(x)$, a
star edge $(\sigma,x)$ of that weight.  If $\alpha_s(x)=0$ (which occurs when
$s=x\in L$), we regard $x$ as an additional copy of the source rather than as
the endpoint of a positive-weight edge.  Call the resulting weighted graph
$A_s$.  Its restriction to $L$ is $Q_L$, and
$d_{A_s}(\sigma,u)=\delta_{L,s}(u)$.

By \eqref{eq:shortcut-properties}, after its first star edge a shortest
$\sigma$--$u$ path can be chosen to have fewer than $4m/k$ further edges.
Consequently every such distance has a shortest realization with at most
\[
  h_Q=\left\lceil\frac{4m}{k}\right\rceil+1
\]
edges.  Apply Nanongkai's scaled light-weight $h_Q$-hop SSSP routine to $A_s$.
The only point not covered verbatim by a physical source in the overlay is the
first star edge.  At scale $j$, landmark $x$ locally forms exactly the rounded
weight that this edge would have in the routine,
\[
  \left\lceil
    \frac{2h_Q\alpha_s(x)}{\eps 2^j}
  \right\rceil,
\]
and schedules the initial label/message that it would receive from $\sigma$.
Infinite labels schedule nothing, and a zero label is present at time zero.
After this initialization every relaxation and every message is an ordinary
operation on an embedded edge of $Q_L$.  Thus the simulated execution is
identical to the execution on $A_s$ after the silent first hop.

The rounding proof for light-weight $h_Q$-hop SSSP gives the one-sided
guarantee \eqref{eq:virtual-root-guarantee}.  The polynomial weight bound
requires only $O(\log n)$ scales and $O(\log n)$-bit labels.  Simulating the
$h_Q/\eps$ overlay rounds through the BFS tree and then pipelining the $m$
answers costs
\[
  \wtO(Dh_Q/\eps+m)
  =\wtO\!\left(D+\frac{mD}{\eps k}+m\right),
\]
as claimed.  This adapts the light-weight SSSP routine of Theorem 3.2 and the
overlay simulation of Lemma 4.6 in Nanongkai's arXiv v2; Theorem 3.10 supplies
the shortcut-diameter bound.  The local injection above accounts explicitly
for the one additional virtual hop.
\end{proof}

The overlay labels reconstruct distances to all network vertices.

\begin{lemma}[Rooted reconstruction]
\label{lem:rooted-reconstruction}
Suppose $L$ is $h$-hitting and that \eqref{eq:hop-estimate} holds for
$L\cup\{s\}$.  Let $Q_L$ and $\widetilde D_{L,s}$ be as above, set
$\widetilde D_{L,s}(s)=0$, and define
\[
  \widetilde d_{L,s}(v)
  =\min_{u\in L\cup\{s\}}
    \bigl(\widetilde D_{L,s}(u)+\widetilde d^h(u,v)\bigr).
  \tag{7}\label{eq:rooted-reconstruction}
\]
Then, simultaneously for all $v\in V$,
\[
  \dist(s,v)\leq\widetilde d_{L,s}(v)
  \leq(1+\eps)^2\dist(s,v).
  \tag{8}\label{eq:rooted-guarantee}
\]
\end{lemma}

\begin{proof}
Every finite bounded-hop estimate is at least the true distance.  Every edge
of $Q_L$ represents an $H_L$-path, so it too is at least the true distance
between its endpoints.  The virtual-root routine returns the weight of an
approximate route in this graph.  Each term minimized in
\eqref{eq:rooted-reconstruction} is therefore the weight of a sequence from
$s$ to $v$, with every segment no shorter than its true endpoint distance.
The triangle inequality proves the lower bound.

Fix $v$ and inspect the canonical shortest path $P_{sv}$.  If it has at most
$h$ edges, choosing $u=s$ in \eqref{eq:rooted-reconstruction} gives an upper
bound of $(1+\eps)\dist(s,v)$.  Otherwise list the landmarks on $P_{sv}$ in
their path order and let $x$ and $y$ be the first and last.  Since every block
of $h$ consecutive vertices is hit, the prefix $s$--$x$, the suffix $y$--$v$,
and the segment between every two consecutive listed landmarks have at most
$h$ edges.  Their bounded-hop estimates form an $H_L$-walk, and hence
\[
  \alpha_s(x)+d_{Q_L}(x,y)
  =\alpha_s(x)+d_{H_L}(x,y)
  \leq(1+\eps)\dist(s,y).
\]
The virtual-root guarantee and the bounded-hop estimate on the final segment
therefore imply
\begin{align*}
  \widetilde d_{L,s}(v)
  &\leq \widetilde D_{L,s}(y)+\widetilde d^h(y,v)\\
  &\leq (1+\eps)^2\dist(s,y)
       +(1+\eps)\dist(y,v)\\
  &\leq (1+\eps)^2\dist(s,v).
\end{align*}
\end{proof}

\section{Separated sampling and trial evaluation}
\label{sec:algorithm}

Fix sufficiently large absolute constants $a$ and $b$.  For each trial
$i\in[n]$, include every vertex in $L_i$ independently with probability
$r/n$, and independently include every vertex in $C_i$ with probability
$c/n$.  We use
\[
  h=\left\lceil\frac{bn\log n}{r}\right\rceil.
  \tag{9}\label{eq:hitting-hop}
\]
For the asymptotic parameter choices below, $a\log n\leq r\leq c\leq n$
and $h\leq n$ once $n$ exceeds an absolute constant.  Smaller instances use
the exact fallback.

\begin{lemma}[Simultaneous sampling event]
\label{lem:sampling}
With high probability, all of the following hold simultaneously:
\begin{enumerate}[label=(\roman*)]
  \item $r/2\leq|L_i|\leq2r$ and $c/2\leq|C_i|\leq2c$ for every trial;
  \item every $L_i$ is $h$-hitting; and
  \item every vertex belongs to between $c/2$ and $2c$ candidate sets.
\end{enumerate}
\end{lemma}

\begin{proof}
The three relevant binomial expectations are $r$, $c$, and $c$,
respectively.  Chernoff bounds, followed by a union bound over the $n$ trials
and $n$ vertices, prove (i) and (iii) when $a$ is sufficiently large.

For (ii), a canonical path has at most $n$ blocks of $h$ consecutive
vertices, so all $n^2$ canonical paths have at most $n^3$ such blocks.  For a
fixed block $B$ and trial $i$,
\[
  \Pr[L_i\cap B=\varnothing]
  \leq (1-r/n)^h
  \leq \exp(-rh/n)
  \leq n^{-b}.
\]
There are $n$ trials, so the union bound is at most $n^{4-b}$.
\end{proof}

Condition on this event for the moment, and write $m_i=|L_i|$.  We choose a
nominal shortcut parameter $k\leq r/16$.  Hence, for all sufficiently large
$n$, $1\leq k<m_i$ in every good trial.  This constant slack is necessary:
choosing $k=r$ would be invalid whenever the realized landmark sample has
fewer than $r$ vertices.

For trial $i$, the network performs the following reversible preprocessing.
It computes the bounded-hop estimates from $L_i\cup C_i$, orders the candidate
identifiers at the leader and pads the list to length $2c$, and constructs the
$k$-shortcut graph on $L_i$.  A dummy list entry has value $0$ for a maximum
and $+\infty$ for a minimum.  Given a genuine candidate $s$, the virtual-root
routine and \eqref{eq:rooted-reconstruction} give
\[
  \widetilde e_i(s)=\max_{v\in V}\widetilde d_{L_i,s}(v).
\]
The maximum is convergecast to the leader.  Define
\[
  F_\Delta(i)=\max_{s\in C_i}\widetilde e_i(s),
  \qquad
  F_R(i)=\min_{s\in C_i}\widetilde e_i(s).
  \tag{10}\label{eq:trial-values}
\]
Trials failing the explicit size tests are assigned the same extreme dummy
values.  This makes the evaluation circuit total even outside the event of
\Cref{lem:sampling}.

\begin{lemma}[One trial]
\label{lem:one-trial}
On the event of \Cref{lem:sampling}, every candidate satisfies
\[
  \ecc(s)\leq\widetilde e_i(s)
  \leq(1+\eps)^2\ecc(s).
  \tag{11}\label{eq:eccentricity-guarantee}
\]
After preprocessing, a candidate value can be evaluated in
\[
  T_{\mathrm{cand}}
  =\wtO\!\left(D+\frac{rD}{k}+r\right)
  \tag{12}\label{eq:candidate-cost}
\]
rounds.  Either trial value in \eqref{eq:trial-values} can be computed by an
inner quantum optimizer, including preprocessing, in
\[
  T_{\mathrm{trial}}
  =\wtO\!\left(
    D+\frac nr+c+rk
    +\sqrt c\left(D+\frac{rD}{k}+r\right)
  \right)
  \tag{13}\label{eq:trial-cost}
\]
rounds.
\end{lemma}

\begin{proof}
The estimate \eqref{eq:eccentricity-guarantee} follows by maximizing
\eqref{eq:rooted-guarantee} over $v$.  Since $m_i=\Theta(r)$,
\Cref{lem:virtual-root} costs the quantity in
\eqref{eq:candidate-cost}; the final convergecast is absorbed by its $D$ term.

The hop-distance computation has $|L_i\cup C_i|=O(c)$ sources and, by
\eqref{eq:hitting-hop}, costs
\[
  \wtO(D+h/\eps+c)=\wtO(D+n/r+c).
\]
Ordering and broadcasting the candidate list costs $O(D+c)$.  Constructing
the shortcut graph costs $\wtO(D+rk)$.  Exact quantum maximum or minimum over
the padded list of $2c$ entries takes
$\wtO(\sqrt c\,T_{\mathrm{cand}})$ rounds by
\Cref{lem:quantum-optimization}.  Adding these contributions proves
\eqref{eq:trial-cost}.
\end{proof}

\subsection{The nested search is a coherent oracle}
\label{sec:coherence}

The outer search queries a trial index in superposition, so a probabilistic
statement about a measured inner optimizer is not enough.  We record the
standard amplification and uncomputation argument in the present setting.

At the start, fix all classical coins: the membership bits defining all
$L_i,C_i$, and the delay and rounding coins for every bounded-hop and overlay
instance that can be queried.  Give each classical subroutine failure
probability at most $n^{-A-4}$ before taking a union bound.  There are only
polynomially many trials, candidates, and subroutines, so with high probability
one event $\mathcal G$ makes every estimate used by every possible query obey
its one-sided guarantee.  Conditioned on $\mathcal G$, preprocessing and
candidate evaluation are fixed reversible circuits controlled by the trial
and candidate registers.

Let $B_i$ consist of reversible trial preprocessing followed by
$q=\Theta(\log n)$ coherent, measurement-free repetitions of the inner
optimizer and a reversible majority computation on their value strings.  All
branches are padded to the same worst-case number of rounds.  The repetitions
use deferred measurement.  For either mode and every $i$, amplification makes
the probability that $B_i$ writes the wrong value at most
\[
  \eta=n^{-K},
\]
where $K$ is an arbitrarily large fixed constant.  On a fresh zeroed work
register, use the circuit
\[
  B_i^\dagger\,
  \operatorname{XOR}_{\mathrm{value}\to\mathrm{answer}}\,
  B_i.
  \tag{14}\label{eq:clean-oracle}
\]
Its action differs in Euclidean norm by at most $2\sqrt\eta$ from the ideal
XOR oracle for $F_\Delta(i)$ or $F_R(i)$.  The same bound holds when $i$ is in
superposition: the controlled-$i$ sectors are orthogonal.

Allocate a fresh zeroed work block for every forward or inverse oracle call of
the outer optimizer.  This avoids making any claim about
\eqref{eq:clean-oracle} on dirty work states.  If the ideal outer algorithm
uses $Q_{\mathrm{out}}$ queries, a hybrid argument changes its final state by
at most $2Q_{\mathrm{out}}\sqrt\eta$.  Here
$Q_{\mathrm{out}}=\wtO(\sqrt{n/c})\leq\wtO(\sqrt n)$, so choosing $K$
sufficiently large makes this error inverse-polynomial.  The repetitions,
reversible sorting, padding, and uncomputation add only polylogarithmic round
factors.  Thus \eqref{eq:trial-cost} is a valid coherent evaluation cost for
the outer search.  This explicit composition is the only place where
unbounded local memory is used materially.

\section{Outer optimization and parameter choices}
\label{sec:outer}

Let $z_\Delta$ be a vertex of eccentricity $\diamw$.  By
\Cref{lem:sampling}, it belongs to at least $c/2$ candidate sets.  In each
corresponding trial,
\[
  F_\Delta(i)\geq\widetilde e_i(z_\Delta)\geq\diamw,
\]
whereas every good trial satisfies
\[
  F_\Delta(i)\leq(1+\eps)^2\diamw.
\]
Thus an unknown upper-threshold search with $\rho=c/(2n)$ finds a value in
this interval.  For the radius, let $z_R$ be a center.  At least $c/2$ trials
satisfy
\[
  F_R(i)\leq\widetilde e_i(z_R)
  \leq(1+\eps)^2\radw,
\]
and every good trial has $F_R(i)\geq\radw$.  The lower-threshold form gives the
same factor $\wtO(\sqrt{n/c})$.  Bad-size trial sentinels cannot meet the
relevant threshold.

Combining this factor with \eqref{eq:trial-cost} gives
\[
  T(r,c,k)=
  \wtO\!\left[
    \sqrt{\frac nc}
    \left(
      D+\frac nr+c+rk
      +\sqrt c\left(D+\frac{rD}{k}+r\right)
    \right)
  \right].
  \tag{15}\label{eq:master-cost}
\]

We now choose the parameters.  Integer parts and the constant factors needed
for $k\leq r/16$ are suppressed in the displays.  They affect neither
\eqref{eq:master-cost} nor the sample event.

If $D\leq n^{1/3}$, take
\[
  r=n^{1/3},\qquad
  c=n^{2/3},\qquad
  k=\Theta(n^{1/6}\sqrt D),\quad k\leq r/16.
  \tag{16}\label{eq:small-D-parameters}
\]
The factor outside the parentheses in \eqref{eq:master-cost} is $n^{1/6}$.
The preprocessing terms obey
\[
  D+n/r+c+rk=O(n^{2/3}),
\]
and
\[
  \sqrt c\left(D+rD/k+r\right)
  =n^{1/3}\left(D+O(n^{1/6}\sqrt D)+n^{1/3}\right)
  =O(n^{2/3}).
\]
Hence $T=\wtO(n^{5/6})$.

If $n^{1/3}<D\leq n^{1/2}$, take
\[
  r=n^{1/3},\qquad
  c=\frac{n^{4/3}}{D^2},\qquad
  k=\Theta(r),\quad k\leq r/16.
  \tag{17}\label{eq:medium-D-parameters}
\]
Now $r\leq c\leq n^{2/3}$ and the outer factor is $n^{-1/6}D$.  Moreover,
\[
  D+n/r+c+rk=O(n^{2/3})
\]
and
\[
  \sqrt c\left(D+rD/k+r\right)
  =\frac{n^{2/3}}D\bigl(O(D)+n^{1/3}\bigr)
  =O(n^{2/3}).
\]
Therefore $T=\wtO(\sqrt nD)$.  For $D>n^{1/2}$, the exact APSP fallback uses
$\wtO(n)$ rounds.  It is also available in every regime, which yields the
minimum with $n$ in \Cref{thm:main}.

Take $\eps=1/\log n$.  On the intersection of the sampling event in
\Cref{lem:sampling} and the classical subroutine event $\mathcal G$ from
\Cref{sec:coherence}, the preceding threshold argument returns values in
\[
  [\diamw,(1+\eps)^2\diamw]
  \quad\text{and}\quad
  [\radw,(1+\eps)^2\radw].
\]
The amplified quantum procedures fail with inverse-polynomial probability by
\Cref{sec:coherence}.  The leader broadcasts the selected value in $O(D)$
rounds, which is absorbed in every displayed bound.  Since
$(1+\eps)^2=1+o(1)$, the approximation and round bounds in
\Cref{thm:main} follow.

\bibliographystyle{alpha}
\bibliography{references}

\end{document}
